\documentclass[10pt,a4paper]{amsart}
\usepackage[margin=19mm]{geometry}
\usepackage{amsmath,amssymb,mathtools}
\usepackage[scr=rsfs,cal=euler]{mathalfa}
\usepackage{xfrac}
\usepackage[unicode,hidelinks,bookmarks=false]{hyperref}
\newcommand{\ie}{i.e.\ }
\newcommand{\cf}{cf.\ }
\newcommand{\resp}{respectively\ }
\newcommand{\Subsection}{\subsection}
\providecommand{\nobreakdash}{\nobreak\hbox{-}}
\setlength{\emergencystretch}{2em}
\multlinegap0pt
\usepackage{amssymb}
\usepackage[shortlabels]{enumitem}
\usepackage{csquotes}
\newtheorem{remark}{Remark}
\newtheorem{proposition}{Proposition}
\title{Calder\'on-Hardy type spaces and the Heisenberg sub-Laplacian}
\author{Pablo Rocha}
\date{Source: arXiv:2505.12163v3 (CC BY 4.0)}
\begin{document}
\maketitle

\noindent\emph{Source and licence.} Calder\'on-Hardy type spaces and the Heisenberg sub-Laplacian. Version arXiv:2505.12163v3 (CC BY 4.0), distributed by arXiv under the Creative Commons Attribution 4.0 International licence, http://creativecommons.org/licenses/by/4.0/, as recorded in the arXiv licence metadata for this exact version and shown on the abstract page https://arxiv.org/abs/2505.12163v3. Full licence notice: \texttt{licenses/arxiv-2505.12163-CC-BY-4.0.txt}.\par\smallskip

\noindent\textit{Editorial scope.} Counted unit: the article's proposition dist (source0.tex lines 619--628). The article's complete original proof of it is delivered with it (source0.tex lines 630--668, one \texttt{proof} environment of 39 lines). No mathematics is written, repaired or re-derived: the statement and the proof are the article's own lines, in the article's own notation, and no retained line is substituted or re-flowed.  Retained uncounted as the source-local context the proof needs: lines 198--200; lines 218--223.  Nothing was omitted from any retained span, so the package carries no omission record.  No retained line contains a citation, so the wrapper carries no bibliography and no bibliography entry.  No retained line contains a figure environment, an image-inclusion command or drawing code, so no image is delivered and the package declares no asset dependency.  Editorial formatting only, in addition: an AMS \texttt{amsart} preamble that restates the article's own declarations and macros used by the retained lines, namely the environments \texttt{remark}, \texttt{proposition} on separate counters, together with the standard packages those lines need.  The retained environments print in this document as Remark 1, Proposition 1 in order of appearance, whereas the article prints the same items with the numbering of its own full document.  The complete native source of the article as distributed in the arXiv e-print for this exact version is bundled unchanged as \texttt{sources/178288/source0.tex}.
\begin{equation} \label{homog dim}
\int_{\mathbb{H}^{n}} f(r \cdot z) \, dz = r^{-Q} \int_{\mathbb{H}^{n}} f(z) \, dz,
\end{equation}

\begin{remark} \label{cambio de centro}
If $f \in L^{1}(\mathbb{H}^{n})$, then for every $\rho$ - ball $B$ and every $z_0 \in \mathbb{H}^{n}$, we have
\[
\int_{B} f(w) \, dw = \int_{z_{0}^{-1} \cdot B} f(z_0 \cdot u) \, du.
\]
\end{remark}

\begin{proposition} \label{Lg dist} Let $g \in L^q_{loc} \cap \mathcal{S}'(\mathbb{H}^n)$ and $f = \mathcal{L} g$ in 
$\mathcal{S}'(\mathbb{H}^n)$. If $\phi \in \mathcal{S}(\mathbb{H}^n)$ and $N > Q+2$, then
\[
(M_{\phi}f)(z) := \sup \left\{ |(f \ast \phi_t)(w)| : \rho(w^{-1} \cdot z) < t, \, 0 < t < \infty \right\}  
\]
\[
\leq C \| \phi \|_{\mathcal{S}(\mathbb{H}^{n}), N}  \,\,\, \eta_{q, 2}(g; \, z)
\]
holds for all $z \in \mathbb{H}^n$.
\end{proposition}

\begin{proof} Let $\rho(w^{-1} \cdot z) < t$, since $f = \mathcal{L} g$ in $\mathcal{S}'(\mathbb{H}^n)$ a computation gives
\[
(f \ast \phi_t)(w) = t^{-2} (g \ast (\mathcal{L} \phi)_t)(w) = t^{-2} \int g(u) (\mathcal{L} \phi)_t (u^{-1} \cdot w) du.
\]
Applying Remark \ref{cambio de centro} and (\ref{homog dim}), we get
\begin{equation} \label{g conv Lphi}
(f \ast \phi_t)(w) = t^{-2} \int g(z \cdot tu) (\mathcal{L} \phi) (u^{-1} \cdot t^{-1}(z^{-1} \cdot w)) du.
\end{equation}
Being $\rho(z^{-1} \cdot w) < t$, a computation gives
\begin{equation} \label{1 mas ro}
1 + \rho(u) \leq 2 \left( 1 + \rho(u^{-1} \cdot t^{-1}(z^{-1} \cdot w)) \right).
\end{equation}
On the other hand, for $N > 2$, we have
\begin{equation} \label{L ineq}
\left|(\mathcal{L} \phi) (u^{-1} \cdot t^{-1}(z^{-1} \cdot w))\right| \left( 1 + \rho(u^{-1} \cdot t^{-1}(z^{-1} \cdot w)) \right)^N \leq 
\| \phi \|_{\mathcal{S}(\mathbb{H}^{n}), N}.
\end{equation}
Now, from (\ref{1 mas ro}) and (\ref{L ineq}), it follows that
\begin{equation} \label{L ineq 2}
\left|(\mathcal{L} \phi) (u^{-1} \cdot t^{-1}(z^{-1} \cdot w))\right| \leq 2^N \| \phi \|_{\mathcal{S}(\mathbb{H}^{n}), N} (1 + \rho(u))^{-N},
\end{equation}
for $\rho(z^{-1} \cdot w) < t$. Then, (\ref{g conv Lphi}), (\ref{L ineq 2}) and (\ref{homog dim}) give
\[
2^{-N} \| \phi \|_{\mathcal{S}(\mathbb{H}^{n}), N}^{-1} \left|(f \ast \phi_t)(w) \right| \leq  t^{-2} \int |g(z \cdot tu)|(1 + \rho(u))^{-N} du.
\]
\[
= t^{-2} t^{-Q} \int |g(z \cdot u)|(1 + \rho(t^{-1} u))^{-N} du
\]
\[
\leq t^{-2} t^{-Q} \int_{\rho(u) < t} |g(z \cdot u)|(1 + \rho(t^{-1} u))^{-N} du 
\]
\[
+ \,\, t^{-2} t^{-Q} \int_{2^j t \leq \rho(u) < 2^{j+1}t} |g(z \cdot u)| \, \rho(t^{-1} u)^{-N} du 
\]
\[
\lesssim \left( 1 + \sum_{j=0}^{\infty} 2^{j(Q+2-N)} \right) \, \eta_{q, 2}(g; z),
\]
for $\rho(z^{-1} \cdot w) < t$. Applying Jensen's inequality and taking $N > Q+2$ in the last inequality the proposition follows.
\end{proof}

\end{document}
